\subsection{Continuing the actual kernel lattice and coefficient action}
\label{subsec:ftb-relative-period-kernel-action}

Use the actual relative exponential \(e_b\) of
Proposition~\ref{prop:ftb-actual-relative-exponential}, on
an open neighborhood \(B\) of \(b_0\), including its
\(\mathrm{HasLifts}\) clause and uniform radius \(r>0\).
In particular \(G/X\) is separated.
All period vectors in this subsection are vectors in
the domain of these same actual maps \(e_b\).

\begin{srlemma}[The source's lattice-continuation criterion]
\label{lem:ftb-kernel-lattice-continuation}
Let \(L\) be finite free over \(\mathbb Z\), let
\(K_b\subset\mathbb C^g\) be additive subgroups, and let
\(P_b:L\to\mathbb C^g\) be additive.
Assume that on one open neighborhood \(B\) of \(b_0\):
\[
\begin{gathered}
 z\in K_b,\ \|z\|<r\ \Longrightarrow z=0,\quad r>0,\qquad
 P_b(L)\subset K_b;\\
 P_{b_0}\text{ is injective},\quad
 P_{b_0}(L)=K_{b_0},\quad
                  \operatorname{span}_{\mathbb R}P_{b_0}(L)=\mathbb C^g.
                                                               \tag{PK1}
\end{gathered}
Assume \(b\mapsto P_b(x)\) is continuous at \(b_0\) for every
\(x\in L\), and the following sequential closedness:
if \(b_j\to b_0\), \(z_j\to z\), \(z_j\in K_{b_j}\), then
\(z\in K_{b_0}\).
Then near \(b_0\),
\[
 P_b(L)=K_b,\qquad P_b\text{ is injective},\qquad
 \operatorname{span}_{\mathbb R}P_b(L)=\mathbb C^g,\qquad
                         P_b(L)\text{ is discrete}.        \tag{PK2}
\]
\end{srlemma}
\begin{proof}
The uniform small-vector condition makes each \(P_b(L)\)
discrete: its intersection with \(B(0,r)\) is \(\{0\}\),
and translate this open singleton through the subgroup.
Choose a \(\mathbb Z\)-basis \(a_1,\ldots,a_s\) of \(L\).
At \(b_0\) its image is a full real lattice, so its rank
is the real dimension \(2g\).
Hence \(s=2g\), and the vectors \(P_{b_0}(a_i)\) are a
real basis.
Continuity and openness of linear independence keep
\(P_b(a_i)\) a real basis near \(b_0\).
They already imply injectivity and real spanning.

Suppose equality of the groups fails arbitrarily near \(b_0\).
Choose \(b_j\to b_0\) and
\(z_j\in K_{b_j}\setminus P_{b_j}(L)\), with the preceding
real-basis property.
Write \(z_j=\sum_i c_{ji}P_{b_j}(a_i)\).
For every fixed integral vector \(n=(n_i)\), put
\[
 z_{j,n}=\sum_i(\operatorname{fract}(c_{ji})-n_i)P_{b_j}(a_i).
                                                               \tag{PK3}
\]
It is \(z_j\) minus the period with coordinates
\(\lfloor c_{ji}\rfloor+n_i\).
Thus \(z_{j,n}\in K_{b_j}\setminus P_{b_j}(L)\).

Pass to a subsequence on which all the fractional parts
converge to \(c_i\in[0,1]\).
Then \(z_{j,0}\to z_*=\sum_i c_iP_{b_0}(a_i)\).
Sequential closedness and (PK1) give \(z_*=P_{b_0}(x)\)
for some \(x\in L\).
Let \(n_i\) be the \(\mathbb Z\)-coordinates of this fixed \(x\).
Now \(z_{j,n}\to0\).
For large \(j\), the uniform radius in (PK1) makes
\(z_{j,n}=0\).
This contradicts the nonmembership assertion following
(PK3), since \(0\in P_{b_j}(L)\).
It proves equality and all of (PK2).
\end{proof}

\begin{srproposition}[Holomorphic continuation of the full actual kernel]
\label{prop:ftb-holomorphic-full-period-kernel}
Let \(L\) be finite free over \(\mathbb Z\), and prescribe
an injective additive \(P_0:L\to\mathbb C^g\) whose image
is the full kernel of \(e_{b_0}\) and spans \(\mathbb C^g\)
over \(\mathbb R\).
After shrinking \(B\), there are additive \(P_b:L\to\mathbb C^g\)
such that \(P_{b_0}=P_0\), every \(P_b(x)\) is holomorphic
in \(b\), and (PK2) holds with \(K_b=\ker e_b\).
\end{srproposition}
\begin{proof}
For a fixed \(\mathbb Z\)-basis \(a_i\), apply the actual
\(\mathrm{HasLifts}\) clause to the identity section, with
the prescribed central lift \(P_0(a_i)\).
It gives analytic vectors \(v_i(b)\), with
\[
              v_i(b_0)=P_0(a_i),\qquad e_b(v_i(b))=1
\]
near \(b_0\).
There are only finitely many basis vectors.
Shrink once to enforce these equations and analyticity
for all of them, and define \(P_b\) by integral linear
extension.
Then \(P_{b_0}=P_0\), each \(P_b(x)\) is holomorphic, and
\(P_b(L)\subset\ker e_b\).

The kernel family has the required sequential closedness.
If \(e_{b_0}(z)\ne1\), the pair \((b_0,z)\) lies in the
open inverse image of the complement of the algebraic
identity section.
The identity section is closed since \(G/X\) is separated.
For actual complex points, factoring through that section
is equivalent to being the identity point in the fibre:
pull back the closed immersion to the reduced
\(\operatorname{Spec}\mathbb C\).
Thus the set \(e_b(z)\ne1\) is open in the actual
regular-evaluation topology.
A convergent sequence of kernel points cannot have such
a limit.
The common small-kernel radius is supplied by (RE9).
Lemma~\ref{lem:ftb-kernel-lattice-continuation} now proves
that these continued basis periods generate the entire
kernel, with all clauses of (PK2).
\end{proof}

\begin{srlemma}[Holomorphic lifts of the actual coefficient action]
\label{lem:ftb-holomorphic-coefficient-action}
Let \(\mathcal O\) be a commutative ring acting by
algebraic group endomorphisms on the given \(G/X\).
There are ring homomorphisms
\(F_b:\mathcal O\to M_g(\mathbb C)\), with all entries
holomorphic in \(b\), for which
\[
              e_b(F_b(a)v)=[a]e_b(v)
                       \quad(a\in\mathcal O,\ v\in\mathbb C^g).
                                                               \tag{PK4}
\]
\end{srlemma}
\begin{proof}
First fix one algebraic group endomorphism \(f\).
At any \(b_1\in B\), apply \(\mathrm{HasLifts}\) on
the parameter space \(B\times\mathbb C^g\) to
\((b,v)\mapsto f(e_b(v))\), with central lift \(0\)
at \((b_1,0)\).
The resulting jointly analytic lift is \(\zeta(b,v)\).
Its derivative in \(v\) at \(0\),
\[
                              A_b=d_v\zeta(b,0),            \tag{PK5}
\]
is holomorphic in \(b\), because derivatives of jointly
analytic maps are analytic.

For a fixed nearby \(b\), the germ \(\zeta_b\) lifts the
additive map \(v\mapsto f(e_b(v))\).
Subtract \(\zeta_b(0)\), which belongs to \(\ker e_b\).
The difference between this germ at \(v+w\) and the sum
at \(v,w\) lies in \(\ker e_b\) and tends to \(0\).
The small-kernel radius makes it locally additive.
The additive-germ argument therefore makes its global
linear lift equal to the derivative (PK5).
It follows that
\[
                                  e_b(A_bv)=f(e_b(v))
\]
for every \(v\).
A linear lift is unique: the difference of two lifts
takes values in the discrete subgroup \(\ker e_b\);
restriction to a real line makes that difference zero.
Thus lifts constructed near different \(b_1\) agree
on overlaps.
They give one holomorphic matrix-valued lift for \(f\).

Apply this to each actual \([a]\), \(a\in\mathcal O\).
Uniqueness of lifts transfers the algebraic identities
\([1]=1\), \([a+a']=[a]+[a']\), and
\([aa']=[a][a']\) into the matrix identities for
\(F_b\).
Here addition of endomorphisms is the pointwise group
product, so additivity of \(e_b\) gives exactly matrix
addition.
This proves that \(F_b\) is a ring homomorphism and
establishes (PK4), including holomorphicity of every entry.
\end{proof}

\begin{srlemma}[The source's holomorphic simultaneous diagonalization]
\label{lem:ftb-holomorphic-diagonalization}
Suppose \(\mathcal O\) is finite as a \(\mathbb Z\)-module
and \(\mathbb C\otimes_{\mathbb Z}\mathcal O\) is reduced.
Let \(F_b\) be a holomorphic family of matrix representations.
Suppose its central value is
\[
                        F_{b_0}(a)=D(a)
                           :=\operatorname{diag}(\tau_i(a))_i
                                                               \tag{PK6}
\]
for ring maps \(\tau_i:\mathcal O\to\mathbb C\).
There is a holomorphic invertible matrix \(h_b\), with
\(h_{b_0}=I\), such that
\[
                               F_b(a)h_b=h_bD(a)             \tag{PK7}
\]
on a neighborhood of \(b_0\).
Repeated characters \(\tau_i\) are allowed.
\end{srlemma}
\begin{proof}
Put \(C=\mathbb C\otimes_{\mathbb Z}\mathcal O\).
It is a finite-dimensional reduced commutative
\(\mathbb C\)-algebra, hence Artinian and semisimple.
For every character \(\tau:\mathcal O\to\mathbb C\),
extend it to \(\widehat\tau:C\to\mathbb C\).
Its kernel is an ideal generated by an idempotent \(f_\tau\).
Put \(e_\tau=1-f_\tau\).
Then
\[
 \widehat\tau(e_\tau)=1,\qquad
              (1\otimes a)e_\tau=\tau(a)e_\tau
                                    \quad(a\in\mathcal O). \tag{PK8}
\]
Indeed \((1\otimes a)-\tau(a)\) lies in the kernel
\((f_\tau)\), and multiplication by \(1-f_\tau\) kills it.

Extend \(F_b\) complex linearly to \(\widehat F_b:C\to M_g(\mathbb C)\).
Every entry \(\widehat F_b(c)_{ij}\) is holomorphic,
by writing \(c\) as a finite sum of simple tensors.
Let the \(j\)-th column of \(h_b\) be
\[
                              \widehat F_b(e_{\tau_j})\,\mathbf e_j.
\]
At \(b_0\), (PK6) makes this column \(\mathbf e_j\),
since its \(j\)-th entry is
\(\widehat\tau_j(e_{\tau_j})=1\).
Thus \(h_{b_0}=I\), even when some of the characters repeat.
Apply \(\widehat F_b\) to (PK8) to obtain (PK7)
column by column.
Holomorphicity and \(h_{b_0}=I\) allow restriction to the
open set \(\det h_b\ne0\).
\end{proof}

\begin{srproposition}[The same kernel with the fixed diagonal action]
\label{prop:ftb-fixed-diagonal-period-continuation}
Assume \(L\) is also an \(\mathcal O\)-module, with
\(\mathcal O\) as in the preceding lemma.
Prescribe the central diagonal action (PK6), and suppose
\[
                     P_0(ax)=D(a)P_0(x).
\]
After shrinking to one ball \(C_0\) about \(b_0\), set
\[
                  P'_b=h_b^{-1}P_b,\qquad e'_b(v)=e_b(h_bv).
                                                               \tag{PK9}
\]
Then \(P'_{b_0}=P_0\), \(e'_{b_0}=e_0\);
the periods and exponentials are holomorphic in their
actual respective variables; and
\[
\begin{gathered}
 \ker e'_b=P'_b(L),\quad P'_b\text{ is injective, discrete
 and real spanning};\\
 P'_b(ax)=D(a)P'_b(x),\qquad
                    e'_b(D(a)v)=[a]e'_b(v).               \tag{PK10}
\end{gathered}
The two vertical tangent-bijection clauses hold for \(e'_b\).
\end{srproposition}
\begin{proof}
Choose a ball where the preceding full-kernel and
diagonalization conclusions hold.
For each \(a,x\), the continuous vector
\[
                         d_b=P_b(ax)-F_b(a)P_b(x)
\]
belongs to \(\ker e_b\), by (PK4), and \(d_{b_0}=0\).
Its norm is either \(0\) or at least \(r\), by the common
small-kernel condition.
On the connected ball, the intermediate value theorem
for this continuous norm forbids a nonzero value:
otherwise the norm would take the value \(r/2\).
Hence \(P_b(ax)=F_b(a)P_b(x)\) for every \(a,x\).
There is no infinite intersection of smaller neighborhoods
here; connectedness proves each equality on the same ball.

Now use \(F_b(a)h_b=h_bD(a)\) and (PK9) to obtain the two
identities (PK10).
Matrix inversion is holomorphic on the nonzero-determinant
open set, so all \(P'_b(x)\) are holomorphic.
Composition by \((b,v)\mapsto(b,h_bv)\) proves joint
regular-evaluation holomorphicity of \(e'_b\).
The complex linear isomorphism \(h_b\) identifies the old
and new kernels, preserving injectivity, full real span
and discreteness.
The same chain rule and its inverse preserve the two
tangent-bijection clauses.
Central equality follows from \(h_{b_0}=I\).
\end{proof}

For the actual host packet \(\mathcal O=\mathcal O_E\),
the reducedness premise in the diagonalization lemma holds.
Indeed \(\mathbb C\otimes_{\mathbb Z}E\) is reduced:
\(\mathbb Q/\mathbb Z\) is a localization and
\(E/\mathbb Q\) is finite separable, so the resulting
formally unramified complex algebra is reduced.
Flatness of \(\mathbb C\) over \(\mathbb Z\) makes
\(\mathbb C\otimes\mathcal O_E\to\mathbb C\otimes E\)
injective; a subring of a reduced ring is reduced.
For the actual \(D'.L\), finiteness over \(\mathbb Z\)
follows from finiteness over \(\mathcal O_E\), and its
integral freeness is the one in the source datum.

\paragraph{Frozen source concordance and remaining transport.}
The files below are under
\srcid{src/Development/AbelianVarietyTheory/Proofs/ModuliProofs/}
at \srcid{alpoge-source-7cbf7a0c}.
The compact fractional-coordinate proof (PK3) is
\srcid{LatContMapContinuous.lean}
(\srcid{699599b2518f18b81e3153c10863eaef229a3063}),
the theorem \srcid{latContinuation}.
The actual lifting, full-kernel equality and final
coordinate change are
\srcid{HolomorphicMapModuliLiftsLocallyToPeriods.lean}
(\srcid{d24831bc0c201893d335bf4e390212761ac34d21}),
notably \srcid{LocLift.exists_periods} and the corresponding
steps in \srcid{LocLift.holomorphic_map_lifts_to_periods}.
The algebraic identity-section closedness, linear-germ
lift uniqueness and norm-gap intermediate-value argument
are in \srcid{ExistsLinearEquivLift.lean}
(\srcid{41cc1a3d8251322b2e8ac7a2c9b0c1878dac3601}).
The actual derivative lift (PK5), its overlap uniqueness
and the ring representation are
\srcid{ExistsHolLinearLift.lean}
(\srcid{600ba2a17fe895a89eb8de8c2fa5ddc5e6784cc5}).
The exact idempotent-column construction (PK8) is
\srcid{DifferentiableOnComplexExtend.lean}
(\srcid{eb1cdb3dc7d703cef6026e1330a10afe231f97e9}).

This proves continuation of the actual full period kernel
and its \(\mathcal O_E\)-action, with the original central
maps unchanged.
The integer polarization form, the Weil relations for all
torsion orders, and the two specified finite-level markings
are transported in the next step; they have not been
inferred merely from (PK10).
